\sloppy
\emergencystretch=5em
\hbadness=10000
\vbadness=10000
\tolerance=2000

\providecolor{codebg}{RGB}{248,248,248}
\providecolor{codeframe}{RGB}{200,200,200}
\providecolor{keyword}{RGB}{0,0,180}
\providecolor{string}{RGB}{163,21,21}
\providecolor{comment}{RGB}{0,128,0}

\lstset{
	backgroundcolor=\color{codebg},
	frame=single,
	rulecolor=\color{codeframe},
	language=Python,
	basicstyle=\ttfamily\scriptsize,
	keywordstyle=\color{keyword}\bfseries,
	stringstyle=\color{string},
	commentstyle=\color{comment}\itshape,
	showstringspaces=false,
	breaklines=true,
	breakatwhitespace=true,
	numbers=left,
	numberstyle=\tiny\color{gray},
	numbersep=8pt,
	tabsize=4,
	literate=
	{ف}{{{\bfseries\itshape ف}}}1
	{ی}{{{\bfseries\itshape ی}}}1
}

\section{نمای کلی و هدف ماژول}

ماژول \lr{\texttt{absolute\_crypto\_validation.py}} یک \emph{سوییت اعتبارسنجی مطلق کریپتوگرافیک} است که نقشی متمایز و مکمل نسبت به سایر اسکریپت‌های آزمون در فریم‌ورک شبیه‌سازی \lr{QKD} ایفا می‌کند. در حالی که ماژول \lr{\texttt{validate\_physics.py}} بر آزمون آماری توزیع‌های منبع (\lr{Poisson} و \lr{Thermal}) متمرکز است و از آزمون \lr{Z-test} دودویی برای راستی‌آزمایی \lr{Gain} و \lr{QBER} استفاده می‌نماید، این ماژول به \emph{صفت کریپتوگرافیک هسته‌ای} پروتکل می‌پردازد: آیا کمیت‌های بنیادین امنیت کلید نظیر \lr{single-photon yield} ($Y_1$)، \lr{single-photon QBER} ($e_1$) و \lr{secure key rate} ($R$) که موتور کامل شبیه‌سازی محاسبه می‌کند، با معادلات تحلیلی بسته‌ای که در ادبیات کلاسیک \lr{decoy-state QKD} (به‌ویژه کار مرجع \lr{Lo-Ma-Chen 2005}) پذیرفته شده‌اند، انطباق عددی دارند؟

هدف اصلی این ماژول، پاسخ به پرسش زیر است: \emph{آیا پیاده‌سازی کامل موتور شبیه‌سازی، از منبع \lr{WCP} و کانال فیبر و آشکارساز \lr{SPAD} و \lr{sifting} پایه‌ها و پروتکل \lr{BB84-Decoy} و اثبات امنیتی \lr{Lim-2014} تا استخراج کلید نهایی، خروجی عددی تولید می‌کند که از نظر کریپتوگرافیک معتبر باشد؟} این پرسش از آنجا اهمیت می‌یابد که در شبیه‌سازی‌های عددی پیچیده، خطاهای ظریف در پیاده‌سازی — مانند اشتباه در فرمول \lr{Bayes} برای استخراج $Y_1$، یا نادیده گرفتن کسر \lr{sifting} $q_{\text{sift}}$ در نرخ کلید، یا اشتباه در نرمال‌سازی \lr{QBER} — می‌توانند منجر به نتایجی شوند که از نظر \lr{QBER} ظاهراً قابل‌قبول‌اند اما از نظر امنیت کریپتوگرافیک \emph{غیرقابل دفاع} هستند. به‌عنوان مثال، اگر استخراج $Y_1$ به‌جای کران پایین \lr{Lo2005} از کران بالا استفاده کند، شبیه‌سازی نرخ کلیدی گزارش می‌دهد که در عمل قابل دستیابی نیست و هر مقاله‌ای که بر این پایه منتشر شود، از نظر رمزنگاری قابل نقض خواهد بود.

ساختار کلی این ماژول شامل سه تابع تحلیلی و یک تابع اجرای آزمون است. تابع \lr{\texttt{H2}} آنتروپی باینری شانون را با برش عددی (\lr{numerical clipping}) پیاده‌سازی می‌کند. تابع \lr{\texttt{channel\_transmittance}} ضریب انتقال کانال فیبر را بر اساس قانون \lr{Beer-Lambert} محاسبه می‌نماید. تابع \lr{\texttt{get\_analytics}} مجموعه‌ای از پنج کمیت تحلیلی کلیدی — $Q_\mu$، $E_\mu$، $Y_1$، $e_1$ و $R_{\text{asymp}}$ — را بر اساس معادلات بسته \lr{decoy-state} محاسبه می‌کند و آن‌ها را در قالب یک دیکشنری برمی‌گرداند. تابع \lr{\texttt{run\_absolute\_test}} آزمون اصلی است که با مقداردهی پیکربندی پیش‌فرض \lr{\texttt{DEFAULT\_CONFIG}} از ماژول \lr{\texttt{main\_optimized}}، اجرای شبیه‌سازی کامل، خواندن خروجی \lr{CSV} و مقایسه با مقادیر تحلیلی، در پنج پارامتر با تلورانس‌های متفاوت، نتیجه نهایی \lr{PASS}/\lr{FAIL} را اعلام می‌کند.

ویژگی برجسته این ماژول، \emph{تفکیک تلورانس‌ها بر اساس حساسیت کریپتوگرافیک} است. تلورانس $25\%$ برای $Q_\mu$ به‌کار رفته است، زیرا این کمیت در شبیه‌سازی مونت‌کارلو با $10^7$ پالس، نویز آماری قابل‌توجهی دارد و انحراف از مقدار تحلیلی طبیعی است. در مقابل، تلورانس $1\%$ برای $Y_1$، $e_1$ و $R_{\text{asymp}}$ به‌کار رفته، چون این کمیت‌ها از ترکیب خطی چندین کمیت به دست می‌آیند و خطای سیستماتیک در هر یک از آن‌ها می‌تواند به خطای بزرگ در نرخ کلید منجر شود. این تلورانس‌های سخت‌گیرانه، تضمین می‌کنند که هرگونه \lr{regression} در پیاده‌سازی موتور — مثلاً پس از \lr{refactor} کلاس اثبات یا تغییر مدل آشکارساز — بلافاصله شناسایی شود.

در سطح عملیاتی، این ماژول به‌عنوان یک \lr{cryptographic gate} عمل می‌کند؛ یعنی پیش از \lr{merge} هر تغییر در شاخه اصلی و پیش از انتشار هر نتیجه شبیه‌سازی در مقالات یا گزارش‌های فنی، این اسکریپت باید با موفقیت اجرا شود. خروجی کد خروج (\lr{exit code}) این اسکریپت نیز به‌گونه‌ای طراحی شده که در \lr{CI/CD pipeline} قابل استفاده باشد: در صورت موفقیت همه آزمون‌ها، کد خروج $0$ و در غیر این صورت، کد خروج $1$ بازگردانده می‌شود. این الگو در مهندسی نرم‌افزار علمی به‌عنوان \lr{cryptographic regression testing} شناخته می‌شود و در پروژه‌های بالغی مانند \lr{Open Quantum Safe} و \lr{Qiskit} به‌صورت استاندارد به کار می‌رود.

نکته مهم دیگر آن است که این ماژول \emph{بدون اعمال هرگونه \lr{patch}} روی موتور اصلی، آن را آزمون می‌کند. به عبارت دیگر، برخلاف برخی اسکریپت‌های آزمون که با \lr{monkey-patching} یا \lr{mock} کردن بخش‌هایی از کد، آن را برای آزمون قابل‌کنترل می‌کنند، این ماژول از همان مسیر اجرای واقعی \lr{\texttt{main\_optimized.run\_and\_save\_csv}} استفاده می‌کند و فقط پیکربندی را override می‌نماید. این رویکرد تضمین می‌کند که آزمون، رفتار واقعی موتور را در شرایط عملیاتی منعکس می‌کند و نه رفتار آزمایشگاهی مصنوعی.

\section{مبانی تئوری و فیزیکی}

\subsection{آنتروپی باینری شانون و نقش آن در امنیت QKD}

آنتروپی باینری شانون، یکی از مفاهیم بنیادین نظریه اطلاعات است که در رمزنگاری کوانتومی نقشی مرکزی ایفا می‌کند. این کمیت که با $H_2(p)$ نشان داده می‌شود، میزان عدم‌قطعیت (\lr{uncertainty}) متغیر تصادفی باینری با توزیع \lr{Bernoulli} را اندازه می‌گیرد و به‌صورت زیر تعریف می‌شود:

\begin{equation}
	H_2(p) = -p \log_2(p) - (1 - p) \log_2(1 - p)
\end{equation}

در این رابطه، $p \in [0, 1]$ احتمال وقوع یکی از دو حالت متغیر باینری است. تابع $H_2(p)$ تقارن نسبت به $p = 0.5$ دارد، در این نقطه به مقدار بیشینه $1$ می‌رسد و در نقاط $p = 0$ و $p = 1$ به صفر میل می‌کند. این رفتار بازتاب‌دهنده این واقعیت استنباطی است که در حالت $p = 0.5$، متغیر تصادفی کاملاً غیرقابل پیش‌بینی است و در حالت‌های $p = 0$ یا $p = 1$ کاملاً قطعی است.

در رمزنگاری کوانتومی، $H_2(p)$ در دو نقطه بحرانی ظاهر می‌شود. نخست، در محاسبه \lr{privacy amplification} که حجم اطلاعات نشت‌کرده به حملهور (\lr{eavesdropper}) را تعیین می‌کند؛ اگر \lr{QBER} سیگنال $E_\mu$ باشد، آنتروپی شرطی شانون $H_2(E_\mu)$ حجم اطلاعات بیت-خطای حملهور را تقریب می‌زند. دوم، در محاسبه \lr{single-photon error rate} $e_1$ که آنتروپی $H_2(e_1)$ حجم اطلاعات حملهور درباره فوتون‌های تک‌فوتونی را تعیین می‌کند و مستقیماً در نرخ کلید امن وارد می‌شود.

پیاده‌سازی تابع $H_2$ در کد، با برش عددی (\lr{numerical clipping}) در نقاط $p = 0$ و $p = 1$ همراه است، زیرا در این نقاط، عبارت $\log_2(p)$ به $-\infty$ میل می‌کند و عبارت $p \log_2(p)$ به حالت تعریف‌نشده $\lim_{p \to 0^+} p \log_2 p = 0$ می‌رسد. در محاسبات عددی کامپیوتری، این حالت می‌تواند منجر به \lr{NaN} یا \lr{exception} شود؛ بنابراین پیاده‌سازی پیش از فراخوانی \lr{\texttt{math.log2}}، شرایط $p \leq 0$ یا $p \geq 1$ را بررسی کرده و در این موارد، مقدار $0$ را برمی‌گرداند. این برش با حد گسسته \lr{L'Hôpital} هم‌خوان است و از نظر ریاضی صحیح است.

\subsection{کاهیدگی فیبر و قانون Beer-Lambert}

انتقال فوتون در فیبر نوری تک‌حالت در طول‌موج مخابراتی $1550$ نانومتر با قانون \lr{Beer-Lambert} مدل می‌شود. این قانون که اصلی‌ترین مدل کاهیدگی در فیبرهای مخابراتی است، رابطه‌ای نمایی بین طول فیبر و ضریب انتقال برقرار می‌سازد. برای یک فیبر به طول $L$ کیلومتر با ضریب کاهیدگی خطی $\alpha_{\text{dB/km}}$، کاهیدگی کل بر حسب دسی‌بل برابر با $\alpha_{\text{dB/km}} \cdot L$ است و تبدیل آن به ضریب انتقال خطی (\lr{power transmittance}) از رابطه زیر پیروی می‌کند:

\begin{equation}
	\eta_{\text{channel}}(L, \alpha_{\text{dB/km}}) = 10^{-\frac{\alpha_{\text{dB/km}} \cdot L}{10}}
\end{equation}

در فیبر تک‌حالت استاندارد در $1550$ نانومتر، $\alpha_{\text{dB/km}} \approx 0.2$ است؛ بنابراین در فاصله $L = 50$ کیلومتر (که نقطه آزمون پیش‌فرض این ماژول است)، ضریب انتقال کانال برابر با $10^{-1} = 0.1$ می‌شود. این کاهش نمایی، عامل اصلی محدودیت برد \lr{QKD} در فیبر است و دلیل آن است که در فواصل بزرگ‌تر از $200$ کیلومتر، حتی با \lr{decoy-state}، نرخ کلید امن به صفر میل می‌کند.

ضریب انتقال کلی سیستم شامل بازده آشکارساز $\eta_{\text{det}}$ نیز می‌شود:

\begin{equation}
	\eta_{\text{total}}(L) = \eta_{\text{det}} \cdot \eta_{\text{channel}}(L, \alpha_{\text{dB/km}})
\end{equation}

در این پیاده‌سازی، $\eta_{\text{det}} = 0.15$ که متناظر با \lr{InGaAs SPAD} در دمای $-50$ درجه سانتی‌گراد است. با این مقادیر، در $L = 50$ کیلومتر، $\eta_{\text{total}} = 0.015$ می‌شود که نشان‌دهنده شرایط عملیاتی واقعی برای \lr{QKD} در فواصل مخابراتی متوسط است.

\subsection{تئوری Decoy-State و معادلات Lo2005}

ایده اصلی \lr{decoy-state QKD} که توسط \lr{Hwang (2003)} پیشنهاد و توسط \lr{Lo-Ma-Chen (2005)} به فرم مدرن آن بسط داده شد، آن است که فرستنده با تغییر تصادفی شدت پالس (\lr{intensity}) در میان چندین سطح (\lr{signal}، \lr{decoy}، \lr{vacuum})، امکان تخمین آماری ویژگی‌های کانال را بدون افشای اینکه کدام پالس حامل کلید است، فراهم می‌سازد. این تکنیک، حمله \lr{Photon-Number Splitting} (\lr{PNS}) را که علیه منبع \lr{WCP} با توزیع \lr{Poisson} قابل اعمال است، خنثی می‌کند و امکان دستیابی به نرخ کلیدی نزدیک به حد امن تک‌فوتونی آرمانی را فراهم می‌سازد.

در ادامه، معادلات بنیادین این تئوری را که همگی در تابع \lr{\texttt{get\_analytics}} پیاده‌سازی شده‌اند، معرفی می‌نماییم.

\subsubsection{Gain سیگنال و QBER سیگنال}

برای منبع \lr{WCP} با توزیع \lr{Poisson} با شدت میانگین $\mu$، \lr{Gain} سیگنال (یعنی احتمال آنکه یک پالس منجر به ثبت کلیک در آشکارساز شود) برابر است با:

\begin{equation}
	Q_\mu^{\text{raw}}(\mu, \eta, Y_0) = 1 - (1 - Y_0) \cdot e^{-\eta \mu}
\end{equation}

در این رابطه، $Y_0$ نرخ \lr{dark count} به ازای هر پالس (احتمال کلیک اشتباه در نبود فوتون واقعی)، $\eta$ ضریب انتقال کلی سیستم و $\mu$ شدت میانگین سیگنال است. عبارت $e^{-\eta \mu}$ احتمال آنکه هیچ فوتونی شناسایی نشود (یعنی یا هیچ فوتونی ارسال نشده، یا فوتون‌های ارسالی از دست رفته‌اند) و عامل $(1 - Y_0)$ احتمال آنکه دارک کانت نیز رخ ندهد را نشان می‌دهد. بنابراین، $1 - (1-Y_0) e^{-\eta \mu}$ احتمال متمم، یعنی احتمال آنکه حداقل یک کلیک ثبت شود، است.

در پروتکل \lr{BB84}، فرستنده و گیرنده هر یک به‌طور تصادفی یکی از دو پایه (\lr{rectilinear} یا \lr{diagonal}) را انتخاب می‌کنند و فقط در حدود نیمی از پالس‌ها (یعنی زمانی که پایه‌ها هم‌خوان باشند) برای استخراج کلید مفید هستند. این کسر به \lr{sifting factor} موسوم است:

\begin{equation}
	q_{\text{sift}} = \frac{1}{2}
\end{equation}

\lr{Gain} سیگنالِ \emph{sifted} شده برابر است با:

\begin{equation}
	Q_\mu^{\text{sifted}} = q_{\text{sift}} \cdot Q_\mu^{\text{raw}}
\end{equation}

\lr{QBER} سیگنال، یعنی کسر کلیک‌های اشتباه در میان کلیک‌های ثبت‌شده، از رابطه زیر به دست می‌آید:

\begin{equation}
	E_\mu = \frac{0.5 \cdot Y_0 + e_d \cdot (1 - e^{-\eta \mu})}{Q_\mu^{\text{raw}}}
\end{equation}

در این رابطه، $e_d$ خطای ذاتی سیستم (\lr{misalignment error}) است که ناشی از عدم‌تقارن پایه‌ها، ناپایداری قطعات اپتیکی و نویز الکترونیکی است. در این پیاده‌سازی، $e_d = 0.01$ به‌کار رفته که در عمل برای سیستم‌های متعادل قابل دستیابی است. صورت این کسر از دو بخش تشکیل شده است: $0.5 \cdot Y_0$ که نشان‌دهنده کلیک‌های اشتباه ناشی از دارک کانت است (چون دارک کانت کاملاً تصادفی است، نصف آن در پایه درست و نصف در پایه اشتباه رخ می‌دهد) و $e_d \cdot (1 - e^{-\eta \mu})$ که نشان‌دهنده کلیک‌های اشتباه ناشی از فوتون‌های واقعی با احتمال خطای ذاتی $e_d$ است.

\subsubsection{Gain دکوی ضعیف}

در پروتکل \lr{decoy-state}، علاوه بر سیگنال با شدت $\mu$، از یک \lr{weak decoy} با شدت $\nu < \mu$ نیز استفاده می‌شود. \lr{Gain} این دکوی ضعیف به‌صورت مشابه محاسبه می‌شود:

\begin{equation}
	Q_\nu(\nu, \eta, Y_0) = 1 - (1 - Y_0) \cdot e^{-\eta \nu}
\end{equation}

این کمیت برای استخراج $Y_1$ (بازده تک‌فوتونی) به‌کار می‌رود. ایده اصلی آن است که با مقایسه $Q_\mu$ و $Q_\nu$ که از دو شدت متفاوت به دست آمده‌اند، می‌توان ویژگی‌های کانال برای فوتون‌های تک‌عددی را استخراج کرد، چون وابستگی $Q_\mu$ به $\mu$ از ترکیب خطی توزیع \lr{Poisson} و ویژگی‌های کانال به دست می‌آید.

\subsubsection{استخراج کران پایین بازده تک‌فوتونی Y1}

مرحله کلیدی در \lr{decoy-state QKD}، استخراج کران پایین برای $Y_1$ (احتمال آنکه یک پالس تک‌فوتونی منجر به کلیک شود) از مشاهدات $Q_\mu$ و $Q_\nu$ است. این کران، که در کار مرجع \lr{Lo-Ma-Chen (2005)} به‌تفصیل استخراج شده، به‌صورت زیر بیان می‌شود:

\begin{equation}
	Y_1^{L} = \frac{\mu^2 e^{-\mu}}{\mu \nu - \mu^2} \left( Q_\nu e^{\nu} - \frac{\nu^2}{\mu^2} Q_\mu^{\text{raw}} e^{\mu} \right)
\end{equation}

این فرمول از ایده زیر نشأت می‌گیرد: $Q_\mu$ را می‌توان به‌صورت جمع وزن‌دار بر روی تعداد فوتون $n$ نوشت:

\begin{equation}
	Q_\mu = \sum_{n=0}^{\infty} Y_n \cdot P(n | \mu) = \sum_{n=0}^{\infty} Y_n \cdot \frac{\mu^n e^{-\mu}}{n!}
\end{equation}

در این رابطه، $Y_n$ بازده کلیک برای پالس $n$-فوتونی است که فرض می‌شود به شدت پالس وابسته نیست (فرض کلیدی \lr{decoy-state}). این فرض بر آن است که حمله حاملهور نمی‌تواند شدت پالس را تشخیص دهد و بنابراین، $Y_n$ برای سیگنال و دکوی یکسان است. با تفریق معادلات $Q_\mu$ و $Q_\nu$ و حذف چندجمله‌ای‌های درجه بالا، می‌توان کران پایین برای $Y_1$ را به دست آورد.

ضریب $\frac{\mu^2 e^{-\mu}}{\mu \nu - \mu^2}$ در فرمول بالا، نرمال‌سازی ریاضی است که از حل سیستم معادلات خطی به دست می‌آید. صورت کران، یعنی $Q_\nu e^{\nu} - \frac{\nu^2}{\mu^2} Q_\mu^{\text{raw}} e^{\mu}$، تفاوت وزن‌دار بین دو مشاهده است که بازده تک‌فوتونی را از نویز چندفوتونی جدا می‌کند. در صورتی که $\mu \approx \nu$ (یعنی شدت سیگنال و دکوی نزدیک به هم باشند)، مخرج $\mu \nu - \mu^2 = \mu(\nu - \mu) \to 0$ میل می‌کند و فرمول به حالت تکین (\lr{singular}) می‌رسد. در این حالت، استخراج $Y_1$ از این روش غیرممکن است و در پیاده‌سازی، مقدار $Y_1^L = 0$ گرفته می‌شود.

\subsubsection{نرخ خطای تک‌فوتونی e1}

نرخ خطای تک‌فوتونی $e_1$، یعنی کسر کلیک‌های اشتباه در میان پالس‌های تک‌فوتونی، از رابطه زیر به دست می‌آید:

\begin{equation}
	e_1 = \frac{0.5 \cdot Y_0 + e_d \cdot \eta}{Y_1}
\end{equation}

این رابطه مشابه $E_\mu$ است، با این تفاوت که به‌جای احتمال کلیک فوتون واقعی $1 - e^{-\eta \mu}$، احتمال کلیک تک‌فوتونی $\eta$ قرار دارد. این جایگزینی برآمده از آن است که برای یک پالس تک‌فوتونی، احتمال شناسایی فوتون دقیقاً برابر با $\eta$ (ضریب انتقال کلی) است، در حالی که برای پالس چندفوتونی، این احتمال ترکیبی از احتمالات کلیک برای هر تعداد فوتون است. در پیاده‌سازی، شرط $Y_1 > 0$ بررسی می‌شود تا از تقسیم بر صفر جلوگیری شود؛ در صورت $Y_1 = 0$، مقدار $e_1 = 0$ گرفته می‌شود.

\subsubsection{بازده تک‌فوتونی و نرخ کلید مجانبی}

\lr{Gain} تک‌فوتونی، یعنی کسری از کلیک‌های کل که از پالس‌های تک‌فوتونی ناشی می‌شوند، از رابطه زیر به دست می‌آید:

\begin{equation}
	Q_1 = \mu \cdot e^{-\mu} \cdot Y_1
\end{equation}

این رابطه از ضرب احتمال آنکه پالس تک‌فوتونی باشد ($P(1 | \mu) = \mu e^{-\mu}$) در بازده تک‌فوتونی $Y_1$ به دست می‌آید. در نهایت، نرخ کلید امن مجانبی (\lr{asymptotic secure key rate}) از فرمول کلاسیک \lr{GLLP} به دست می‌آید:

\begin{equation}
	R_{\text{asymp}} = q_{\text{sift}} \cdot \left[ Q_1 \cdot (1 - H_2(e_1)) - Q_\mu^{\text{raw}} \cdot f_{\text{EC}} \cdot H_2(E_\mu) \right]
\end{equation}

این فرمول، که در ادبیات \lr{QKD} به‌عنوان \lr{GLLP key rate formula} شناخته می‌شود، از دو بخش تشکیل شده است. بخش اول $Q_1 \cdot (1 - H_2(e_1))$ نشان‌دهنده حجم اطلاعات مشترک بین فرستنده و گیرنده از پالس‌های تک‌فوتونی پس از \lr{privacy amplification} است؛ عبارت $1 - H_2(e_1)$ ضریب \lr{privacy amplification} است که حجم اطلاعات نشت‌کرده به حاملهور را کم می‌کند. بخش دوم $Q_\mu^{\text{raw}} \cdot f_{\text{EC}} \cdot H_2(E_\mu)$ نشان‌دهنده حجم اطلاعات فاش‌شده در مرحله \lr{error correction} است؛ $f_{\text{EC}} \geq 1$ بازدهی \lr{error correction} است که در عمل $f_{\text{EC}} \approx 1.1$ برای کد \lr{LDPC} بهینه است. در این پیاده‌سازی، $f_{\text{EC}} = 1.1$ به‌کار رفته است.

تفاوت این فرمول با حالت \lr{finite-key} (مثل اثبات \lr{Lim-2014}) آن است که در اینجا از حد مجانبی $N \to \infty$ استفاده می‌شود و جریمه‌های آماری نظیر \lr{Hoeffding bound} و \lr{composable security} اعمال نمی‌گردند. به همین دلیل، $R_{\text{asymp}}$ همیشه بزرگ‌تر یا مساوی با نرخ کلید متناهی است و در عمل به‌عنوان \emph{کران بالای} نرخ کلید قابل دستیابی عمل می‌کند. با این حال، تطابق موتور شبیه‌سازی با این کران بالا، شرط لازم — هرچند کافی نیست — برای صحت پیاده‌سازی است.

\subsection{تلورانس اعتبارسنجی و تفسیر کریپتوگرافیک}

انتخاب تلورانس‌های متفاوت برای پنج پارامتر، بازتاب‌دهنده درک عمیق از فیزیک و آمار شبیه‌سازی است. تلورانس $25\%$ برای $Q_\mu$ به این دلیل انتخاب شده که در شبیه‌سازی مونت‌کارلو با $N = 10^7$ پالس، انحراف معیار $Q_\mu$ تقریباً برابر با $\sqrt{Q_\mu (1 - Q_\mu) / N} \approx 10^{-4}$ است؛ با این حال، چون دارک کانت $Y_0 = 3 \times 10^{-6}$ بسیار کوچک است، نویز نسبی می‌تواند تا چندین درصد باشد. تلورانس $5\%$ برای $E_\mu$ سخت‌گیرانه‌تر است، چون $E_\mu$ از نسبت دو کمیت به دست می‌آید و واریانس نسبی آن کمتر است.

تلورانس $1\%$ برای $Y_1$، $e_1$ و $R_{\text{asymp}}$ بسیار سخت‌گیرانه است و دلیل آن آن است که این کمیت‌ها در فرمول کلید نهایی به‌صورت ضرب و تفریق وارد می‌شوند و خطای کوچک در هر یک می‌تواند به خطای بزرگ در $R$ منجر شود. به‌ویژه، $R_{\text{asymp}}$ از تفریق دو کمیت نزدیک به هم به دست می‌آید و حساس به خطای سیستماتیک است. در صورتی که تلورانس $1\%$ برای $R_{\text{asymp}}$ رعایت شود، می‌توان با اطمینان گفت که پیاده‌سازی موتور، از نظر کریپتوگرافیک معتبر است.

\section{تحلیل معماری و ساختار داده‌ها}

\subsection{ساختار کلی ماژول و الگوی Module-Level Functions}

این ماژول از الگوی \lr{module-level functions} پیروی می‌کند، یعنی به‌جای تعریف کلاس‌ها، از توابع مستقل در سطح ماژول استفاده می‌نماید. این انتخاب طراحی، بازتاب‌دهنده ماهیت \lr{procedural} این ماژول است: توابع تحلیلی (\lr{\texttt{H2}}، \lr{\texttt{channel\_transmittance}}، \lr{\texttt{get\_analytics}}) توابع ریاضی خالص (\lr{pure functions}) هستند که ورودی‌هایشان را به خروجی‌هایشان نگاشت می‌کنند و هیچ \lr{state} داخلی ندارند. این الگو مزایای زیر را به همراه دارد: تست‌پذیری آسان (بدون نیاز به \lr{setup} و \lr{teardown}، هر تابع مستقل قابل آزمون است)، قابلیت ترکیب (خروجی یک تابع مستقیماً به ورودی تابع بعدی می‌رود) و قابلیت خواندن بالا (جریان کنترل از بالا به پایین واضح است).

تنها تابع دارای \lr{side-effect}، \lr{\texttt{run\_absolute\_test}} است که فایل \lr{CSV} می‌نویسد و گزارش چاپ می‌کند. این جداسازی بین توابع خالص و توابع دارای \lr{side-effect}، یکی از اصول مهندسی نرم‌افزار بالغ است و در الگوهای طراحی نظیر \lr{functional core, imperative shell} تبیین می‌شود. در این الگو، هسته محاسباتی (\lr{functional core}) از توابع خالص تشکیل می‌شود و لایه بیرونی (\lr{imperative shell}) مسئول \lr{I/O} و مدیریت \lr{side-effect} است.

\subsection{تحلیل تابع H2}

تابع \lr{\texttt{H2(p)}} یک تابع ریاضی خالص است که آنتروپی باینری شانون را محاسبه می‌نماید. این تابع با بررسی شرایط $p \leq 0$ یا $p \geq 1$، از خطای \lr{math domain error} در فراخوانی \lr{\texttt{math.log2}} جلوگیری می‌کند. در صورتی که $p$ خارج از بازه $(0, 1)$ باشد، مقدار $0$ بازگردانده می‌شود که با حد ریاضی هم‌خوان است. این پیاده‌سازی، الگوی \lr{defensive programming} را به نمایش می‌گذارد: ورودی‌ها پیش از استفاده در محاسبات، اعتبارسنجی می‌شوند و در صورت غیرمعتبر بودن، مقدار پیش‌فرض امن بازگردانده می‌شود.

\begin{LTR}
	\begin{lstlisting}[language=Python]
		def H2(p):
		if p <= 0 or p >= 1: return 0.0
		return -p * math.log2(p) - (1 - p) * math.log2(1 - p)
	\end{lstlisting}
\end{LTR}

در خط اول، شرط $p \leq 0$ یا $p \geq 1$ بررسی می‌شود. این شرط شامل نقاط انتهایی $p = 0$ و $p = 1$ نیز می‌شود که در آن‌ها آنتروپی به‌صورت تحلیلی صفر است. در خط دوم، فرمول اصلی $-p \log_2 p - (1-p) \log_2(1-p)$ پیاده‌سازی می‌شود. چون این خط فقط در صورتی اجرا می‌شود که $p \in (0, 1)$، فراخوانی \lr{\texttt{math.log2}} همیشه با ورودی مثبت انجام می‌شود و خطای \lr{math domain error} رخ نمی‌دهد.

\subsection{تحلیل تابع channel\_transmittance}

این تابع، ضریب انتقال کانال فیبر را بر اساس قانون \lr{Beer-Lambert} محاسبه می‌کند. این تابع نیز یک تابع خالص است که دو ورودی $L$ (طول فیبر به کیلومتر) و $\alpha_{\text{dB/km}}$ (ضریب کاهیدگی به \lr{dB/km}) را می‌گیرد و ضریب انتقال خطی را برمی‌گرداند.

\begin{LTR}
	\begin{lstlisting}[language=Python]
		def channel_transmittance(L, alpha_dB_km):
		return 10 ** (-alpha_dB_km * L / 10)
	\end{lstlisting}
\end{LTR}

این پیاده‌سازی ساده، مستقیماً فرمول $\eta_{\text{channel}} = 10^{-\alpha_{\text{dB/km}} \cdot L / 10}$ را پیاده می‌سازد. انتخاب نمادگذاری $10^{**}$ به‌جای \lr{\texttt{math.pow}} یا \lr{\texttt{exp}}، عمدی است: در پایتون، عملگر \lr{**} برای توان‌دهی عددی به‌کار می‌رود و خوانایی بالاتری دارد. در این فرمول، علامت منفی در نماد $-\alpha_{\text{dB/km}} \cdot L / 10$ بازتاب‌دهنده آن است که کاهیدگی باعث کاهش انتقال می‌شود و بنابراین، ضریب انتقال با افزایش طول فیبر، به‌صورت نمایی کاهش می‌یابد.

\subsection{تحلیل تابع get\_analytics و ساختار خروجی}

تابع \lr{\texttt{get\_analytics}} اصلی‌ترین تابع تحلیلی این ماژول است که پنج کمیت کلیدی کریپتوگرافیک را محاسبه می‌کند. این تابع هفت پارامتر ورودی دارد: $\mu$ (شدت سیگنال)، $\nu$ (شدت دکوی ضعیف)، $\eta$ (ضریب انتقال کلی)، $Y_0$ (نرخ دارک کانت)، $e_d$ (خطای ذاتی)، $f_{\text{EC}}$ (بازدهی \lr{error correction}) و $q_{\text{sift}}$ (کسر \lr{sifting}) که مقدار پیش‌فرض $0.5$ دارد. خروجی این تابع یک دیکشنری است که کلیدهای \lr{\texttt{Q\_mu}}، \lr{\texttt{E\_mu}}، \lr{\texttt{Y1}}، \lr{\texttt{e1}} و \lr{\texttt{R\_asymp}} را دارد.

\begin{LTR}
	\begin{lstlisting}[language=Python]
		def get_analytics(mu, nu, eta, Y0, ed, f_ec, q_sift=0.5):
		# 1. پارامترهای سیگنال (mu)
		Q_mu_raw = 1 - (1 - Y0) * math.exp(-eta * mu)
		Q_mu_sifted = Q_mu_raw * q_sift
		E_mu = (0.5 * Y0 + ed * (1 - math.exp(-eta * mu))) / Q_mu_raw
		
		# 2. پارامترهای فریب ضعیف (nu)
		Q_nu = 1 - (1 - Y0) * math.exp(-eta * nu)
	\end{lstlisting}
\end{LTR}

در خطوط ابتدایی، پارامترهای سیگنال محاسبه می‌شوند. $Q_\mu^{\text{raw}}$ از فرمول بسته محاسبه می‌شود؛ $Q_\mu^{\text{sifted}}$ با ضرب در $q_{\text{sift}}$ به دست می‌آید؛ $E_\mu$ از نسبت خطای کل به \lr{Gain} کل محاسبه می‌شود. در ادامه، $Q_\nu$ برای دکوی ضعیف به‌صورت مشابه محاسبه می‌گردد. این مقادیر، پایه محاسبات بعدی هستند.

\begin{LTR}
	\begin{lstlisting}[language=Python]
		# 3. استخراج کران پایین تک‌فوتونیتی (Lo2005)
		if abs(mu - nu) < 1e-9:
		Y1_L = 0
		else:
		term1 = Q_nu * math.exp(nu)
		term2 = (nu**2 / mu**2) * Q_mu_raw * math.exp(mu)
		factor = (mu**2 * math.exp(-mu)) / (mu * nu - mu**2)
		Y1_L = factor * (term1 - term2)
		
		Y1 = max(Y1_L, 0)
	\end{lstlisting}
\end{LTR}

در این بخش، کران پایین $Y_1^L$ محاسبه می‌شود. شرط $\left| \mu - \nu \right| < 10^{-9}$ حالت تکین را شناسایی می‌کند؛ در این حالت، فرمول به‌صورت تحلیلی نامحدود می‌شود و مقدار $Y_1^L = 0$ گرفته می‌شود. در غیر این صورت، فرمول اصلی \lr{Lo2005} پیاده می‌شود: \lr{\texttt{term1}} و \lr{\texttt{term2}} دو جمله اصلی صورت کسر هستند و \lr{\texttt{factor}} ضریب نرمال‌سازی است. در نهایت، \lr{\texttt{max(Y1\_L, 0)}} تضمین می‌کند که $Y_1$ منفی نشود؛ این برش عددی ضروری است، چون در برخی موارد (به‌ویژه وقتی $Q_\nu$ به‌دلیل نویز آماری کوچک‌تر از حد تئوریک باشد)، $Y_1^L$ می‌تواند منفی شود که فیزیکاً بی‌معناست.

\begin{LTR}
	\begin{lstlisting}[language=Python]
		# 4. محاسبه e1
		e1 = (0.5 * Y0 + ed * eta) / Y1 if Y1 > 0 else 0.0
		
		# 5. محاسبه نرخ کلید مجانبی
		Q1 = mu * math.exp(-mu) * Y1
		R_asymp = q_sift * (Q1 * (1 - H2(e1)) - Q_mu_raw * f_ec * H2(E_mu))
		R_asymp = max(R_asymp, 0.0)
	\end{lstlisting}
\end{LTR}

در این بخش، $e_1$، $Q_1$ و $R_{\text{asymp}}$ محاسبه می‌شوند. شرط $Y_1 > 0$ از تقسیم بر صفر جلوگیری می‌کند؛ در صورت $Y_1 = 0$، $e_1 = 0$ گرفته می‌شود. $Q_1$ از ضرب $\mu e^{-\mu}$ (احتمال پالس تک‌فوتونی) در $Y_1$ (بازده تک‌فوتونی) به دست می‌آید. $R_{\text{asymp}}$ از فرمول \lr{GLLP} محاسبه می‌شود و در نهایت، \lr{\texttt{max(R\_asymp, 0.0)}} تضمین می‌کند که نرخ کلید منفی نشود؛ این برش بازتاب‌دهنده آن است که در رژیم‌های با \lr{QBER} بالا، نرخ کلید به‌صورت تحلیلی منفی می‌شود که معنای فیزیکی آن این است که هیچ کلید امنی قابل استخراج نیست و در عمل، مقدار صفر گزارش می‌شود.

\begin{LTR}
	\begin{lstlisting}[language=Python]
		return {
			"Q_mu": Q_mu_sifted, "E_mu": E_mu, "Y1": Y1, "e1": e1, "R_asymp": R_asymp
		}
	\end{lstlisting}
\end{LTR}

خروجی تابع به‌صورت یک دیکشنری با پنج کلید است. این انتخاب طراحی، به دلیل خوانایی بالا و انعطاف‌پذیری در افزودن کمیت‌های جدید در آینده است. در صورتی که در نسخه‌های بعدی، کمیت‌های جدیدی نظیر $Y_0^L$ (کران پایین دارک کانت) یا $e_1^U$ (کران بالا خطای تک‌فوتونی) اضافه شوند، فقط یک کلید جدید به دیکشنری اضافه می‌شود و کدهای فراخوان نیاز به تغییر ندارند.

\subsection{تحلیل تابع run\_absolute\_test و معماری پیکربندی}

تابع \lr{\texttt{run\_absolute\_test}} اصلی‌ترین تابع آزمون است که کل جریان کار را هماهنگ می‌کند. این تابع چهار پارامتر ورودی دارد: $L$ (فاصله به کیلومتر)، $\mu$ (شدت سیگنال)، $\nu$ (شدت دکوی) و $f_{\text{EC}}$ (بازدهی \lr{error correction}). در ابتدا، ماژول \lr{\texttt{main\_optimized}} به‌صورت \lr{lazy} وارد می‌شود:

\begin{LTR}
	\begin{lstlisting}[language=Python]
		def run_absolute_test(dist_km, mu, nu, f_ec):
		try:
		import main_optimized as m
		except ImportError:
		print("Error: main_optimized not found.")
		return False
	\end{lstlisting}
\end{LTR}

این الگوی \lr{lazy import} دو مزیت دارد: نخست، در صورتی که ماژول به‌عنوان کتابخانه وارد شود (نه به‌عنوان اسکریپت اجرا شود)، وارد کردن \lr{\texttt{main\_optimized}} به تأخیر می‌افتد و زمان بارگذاری ماژول کاهش می‌یابد. دوم، در صورتی که \lr{\texttt{main\_optimized}} در دسترس نباشد (مثلاً در محیط \lr{CI} با وابستگی‌های ناقص)، خطا به‌صورت کنترل‌شده مدیریت می‌شود و تابع مقدار \lr{False} برمی‌گرداند.

سپس پیکربندی پیش‌فرض با \lr{\texttt{copy.deepcopy}} کپی می‌شود:

\begin{LTR}
	\begin{lstlisting}[language=Python]
		cfg = copy.deepcopy(m.DEFAULT_CONFIG)
		
		cfg["simulation"]["apply_statistical_noise"] = True
		cfg["simulation"]["distance_start_km"] = dist_km
		cfg["simulation"]["distance_stop_km"] = dist_km
		cfg["simulation"]["distance_points"] = 1
		cfg["simulation"]["min_pulses_log"] = 7
		cfg["simulation"]["max_pulses_log"] = 7
	\end{lstlisting}
\end{LTR}

استفاده از \lr{\texttt{deepcopy}} به‌جای انتساب مستقیم، حیاتی است: \lr{\texttt{DEFAULT\_CONFIG}} یک دیکشنری تودرتوی است و در صورت انتساب مستقیم، تغییرات روی کپی اعمال می‌شود و \lr{\texttt{DEFAULT\_CONFIG}} اصلی نیز تغییر می‌کند. این می‌تواند باعث \lr{regression} در آزمون‌های بعدی شود. \lr{\texttt{deepcopy}} یک کپی عمیق و مستقل ایجاد می‌کند که تغییرات روی آن، روی شیء اصلی اثر نمی‌گذارد. در ادامه، پیکربندی شبیه‌سازی تنظیم می‌شود: فقط یک نقطه فاصله (یعنی $L = $ \lr{\texttt{dist\_km}}) آزمون می‌شود، $N = 10^7$ پالس استفاده می‌شود و نویز آماری فعال است.

\begin{LTR}
	\begin{lstlisting}[language=Python]
		cfg["source"]["intended_proof"] = "LIM_2014"
		cfg["protocol_params"]["protocol"] = "BB84_DECOY"
		cfg["protocol_runtime"]["protocol_class"] = "BB84DecoyProtocol"
		cfg["source"]["intensity_jitter"] = 0.0
		
		cfg["source"]["pulses"] = {
			"signal": {"mu": mu, "prob": 0.5},
			"decoy": {"mu": nu, "prob": 0.25},
			"vacuum": {"mu": 0.0, "prob": 0.25}
		}
	\end{lstlisting}
\end{LTR}

در این بخش، پروتکل و اثبات امنیتی تنظیم می‌شوند. \lr{\texttt{LIM\_2014}} به‌عنوان اثبات امنیتی انتخاب می‌شود، زیرا این اثبات، نسخه متناهی-\lr{key} \lr{decoy-state BB84} است که در فریم‌ورک پیاده‌سازی شده. با این حال، چون فرمول‌های تحلیلی در تابع \lr{\texttt{get\_analytics}} حد مجانبی را محاسبه می‌کنند، آزمون در واقع کران بالای نرخ کلید را راستی‌آزمایی می‌کند. این یک انتخاب طراحی هوشمندانه است: اگر موتور شبیه‌سازی نتواند به حد مجانبی نزدیک شود، قطعاً نمی‌تواند نرخ کلید متناهی صحیح را تولید کند. در ادامه، توزیع پالس‌ها تنظیم می‌شود: $50\%$ سیگنال، $25\%$ دکوی و $25\%$ واکوئم. این توزیع، توزیع کلاسیک \lr{decoy-state} است که در کار \lr{Lim-2014} نیز به‌کار رفته. \lr{\texttt{intensity\_jitter = 0.0}} تضمین می‌کند که شدت پالس‌ها دقیقاً برابر با مقدار تنظیم‌شده باشد و نویز شدت اعمال نشود؛ این برای آزمون ضروری است، چون فرمول‌های تحلیلی فرض می‌کنند که شدت پالس‌ها دقیقاً $\mu$ و $\nu$ هستند.

\begin{LTR}
	\begin{lstlisting}[language=Python]
		eta_det = 0.15
		Y0 = 3e-6
		ed = 0.01
		cfg["detector"]["det_eff_d0"] = eta_det
		cfg["detector"]["det_eff_d1"] = eta_det
		cfg["detector"]["dark_rate"] = Y0
		cfg["detector"]["dark_rate_d1"] = Y0
		cfg["detector"]["qber_intrinsic"] = ed
		cfg["detector"]["misalignment"] = 0.0
		cfg["detector"]["dead_time_ns"] = 0.0
	\end{lstlisting}
\end{LTR}

پارامترهای آشکارساز تنظیم می‌شوند. $\eta_{\text{det}} = 0.15$، $Y_0 = 3 \times 10^{-6}$ و $e_d = 0.01$ انتخاب شده‌اند. این مقادیر با مقادیر مرجع در ادبیات \lr{QKD} هم‌خوان هستند. \lr{\texttt{misalignment = 0.0}} و \lr{\texttt{dead\_time\_ns = 0.0}} تنظیم شده‌اند تا اثرات اضافی حذف شوند و آزمون فقط روی فیزیک اصلی متمرکز باشد. این رویکرد \lr{minimal configuration}، تضمین می‌کند که هرگونه انحراف از فرمول تحلیلی، ناشی از باگ در موتور باشد نه از اثرات فیزیکی اضافی.

\subsection{الگوی Comparison و Tolerance-Based Validation}

قلب تابع \lr{\texttt{run\_absolute\_test}}، تابع داخلی \lr{\texttt{check}} است که مقادیر شبیه‌سازی را با مقادیر تحلیلی مقایسه می‌کند:

\begin{LTR}
	\begin{lstlisting}[language=Python]
		def check(name, fw_val, th_val, tol):
		if th_val == 0: 
		err = abs(fw_val - th_val)
		passed = err < 1e-9
		else:
		err = abs(fw_val - th_val) / abs(th_val) * 100
		passed = err <= tol
		print(f"{name:<15} | {fw_val:<15.6e} | {th_val:<15.6e} | {err:<10.4f} | {tol:<5} | {'PASS' if passed else 'FAIL'}")
		return passed
	\end{lstlisting}
\end{LTR}

این تابع از الگوی \lr{relative tolerance} پیروی می‌کند: در صورتی که مقدار تحلیلی غیرصفر باشد، خطای نسبی به درصد محاسبه می‌شود و با تلورانس مقایسه می‌گردد. در صورتی که مقدار تحیلی صفر باشد، خطای مطلق با آستانه $10^{-9}$ مقایسه می‌شود. این دوگانه‌سازی ضروری است، چون در صورتی که مقدار تحلیلی صفر باشد، خطای نسبی تعریف‌نشده (\lr{division by zero}) است. آستانه $10^{-9}$ برای حالت صفر، بسیار سخت‌گیرانه است و تضمین می‌کند که در این حالت، فقط مقادیر عملاً صفر قابل قبول هستند.

خروجی این تابع، یک خط فرمت‌شده با ستون‌های نام پارامتر، مقدار شبیه‌سازی، مقدار تحلیلی، خطا، تلورانس و نتیجه \lr{PASS}/\lr{FAIL} است. این قالب، خوانایی بالایی دارد و امکان تحلیل سریع نتایج را فراهم می‌سازد. در نسخه اصلی کد، از ایموجی \lr{✅} و \lr{❌} برای نمایش نتیجه استفاده شده، اما در پیاده‌سازی مستند، برای سازگاری با خروجی کنسول استاندارد، می‌توان از \lr{PASS} و \lr{FAIL} استفاده کرد.

\subsection{ساختار پنج‌آزمونی و خروجی نهایی}

پنج آزمون به ترتیب با تلورانس‌های $25\%$، $5\%$، $1\%$، $1\%$ و $1\%$ اجرا می‌شوند:

\begin{LTR}
	\begin{lstlisting}[language=Python]
		p1 = check("Q_mu", fw_Q_mu, th["Q_mu"], 25.0)      
		p2 = check("E_mu", fw_E_mu, th["E_mu"], 5.0)
		p3 = check("Y1 (Extraction)", fw_Y1, th["Y1"], 1.0)  
		p4 = check("e1 (Extraction)", fw_e1, th["e1"], 1.0)
		p5 = check("SKR (Absolute)", fw_R, th["R_asymp"], 1.0) 
		
		all_passed = all([p1, p2, p3, p4, p5])
		if all_passed:
		print("100% CRYPTOGRAPHIC VALIDATION PASSED!")
		else:
		print("MATH/CRYPTO VALIDATION FAILED. CHECK FRAMEWORK ENGINE.")
		
		return all_passed
	\end{lstlisting}
\end{LTR}

در این بخش، پنج آزمون به ترتیب اجرا می‌شوند. تابع \lr{\texttt{all}} بررسی می‌کند که آیا همه نتایج \lr{True} هستند یا خیر. در صورت موفقیت همه، پیام \lr{100\% CRYPTOGRAPHIC VALIDATION PASSED!} چاپ می‌شود و در غیر این صورت، پیام خطا چاپ می‌گردد. خروجی تابع، یک مقدار بولین است که در \lr{\texttt{\_\_main\_\_}} برای تعیین کد خروج استفاده می‌شود. این الگو، تضمین می‌کند که در \lr{CI/CD pipeline}، کد خروج به‌درستی تنظیم شود و در صورت شکست، \lr{pipeline} متوقف گردد.

\subsection{تحلیل نقطه ورود \_\_main\_\_}

\begin{LTR}
	\begin{lstlisting}[language=Python]
		if __name__ == "__main__":
		print("=" * 80)
		print("ABSOLUTE CRYPTOGRAPHIC & MATHEMATICAL VALIDATION SUITE")
		print("Testing Native Framework Engine (No Patching) vs Analytical Ground Truth")
		print("=" * 80)
		
		is_valid = run_absolute_test(dist_km=50.0, mu=0.5, nu=0.1, f_ec=1.1)
		sys.exit(0 if is_valid else 1)
	\end{lstlisting}
\end{LTR}

نقطه ورود \lr{\texttt{\_\_main\_\_}} یک سرآیند تزئینی چاپ می‌کند و سپس تابع \lr{\texttt{run\_absolute\_test}} را با پارامترهای پیش‌فرض $L = 50$ کیلومتر، $\mu = 0.5$، $\nu = 0.1$ و $f_{\text{EC}} = 1.1$ فراخوانی می‌کند. این پارامترها، نقطه عملیاتی متعارف برای \lr{QKD} در فیبر نوری هستند و انتخاب آن‌ها بر اساس کار مرجع \lr{Lim-2014} است. در نهایت، \lr{\texttt{sys.exit}} با کد خروج $0$ (در صورت موفقیت) یا $1$ (در صورت شکست) فراخوانی می‌شود. این الگو، استاندارد \lr{POSIX} برای اسکریپت‌های \lr{CI/CD} است و تضمین می‌کند که \lr{shell} بتواند کد خروج را به‌درستی تفسیر کند.

\section{پیاده‌سازی ریاضی و الگوریتمی}

\subsection{پیاده‌سازی آنتروپی باینری با برش عددی}

پیاده‌سازی تابع $H_2$ در این ماژول، نمونه‌ای از \lr{defensive programming} در محاسبات عددی است. فرمول تحلیلی $-p \log_2 p - (1-p) \log_2(1-p)$ در نقاط $p = 0$ و $p = 1$ به‌صورت تحلیلی به $0$ میل می‌کند، اما در محاسبات عددی، فراخوانی \lr{\texttt{math.log2(0)}} منجر به \lr{ValueError} می‌شود. برای حل این مشکل، شرط $p \leq 0$ یا $p \geq 1$ پیش از فراخوانی \lr{\texttt{math.log2}} بررسی می‌شود و در این موارد، مقدار $0$ بازگردانده می‌شود.

این برش عددی با حد گسسته \lr{L'Hôpital} هم‌خوان است:

\begin{equation}
	\lim_{p \to 0^+} p \log_2 p = \lim_{p \to 0^+} \frac{\log_2 p}{1/p} = \lim_{p \to 0^+} \frac{1/(p \ln 2)}{-1/p^2} = \lim_{p \to 0^+} \frac{-p}{\ln 2} = 0
\end{equation}

بنابراین، پیاده‌سازی عددی با حد تحلیلی هم‌خوان است و از نظر ریاضی صحیح است. این الگو در بسیاری از توابع ریاضی در علوم کامپیوتر به‌کار می‌رود و به‌عنوان \lr{numerical stability via boundary clipping} شناخته می‌شود.

\subsection{پیاده‌سازی فرمول Lo2005 برای Y1}

پیاده‌سازی فرمول $Y_1^L$ در تابع \lr{\texttt{get\_analytics}} چندین نکته مهم دارد. نخست، شرط $\left| \mu - \nu \right| < 10^{-9}$ حالت تکین را شناسایی می‌کند. این شرط بازتاب‌دهنده آن است که در فرمول $Y_1^L$، مخرج $\mu \nu - \mu^2 = \mu(\nu - \mu)$ در صورتی که $\nu \to \mu$ باشد، به صفر میل می‌کند و فرمول به‌صورت تحلیلی نامحدود می‌شود. در عمل، این حالت رخ نمی‌دهد (چون $\nu$ و $\mu$ معمولاً بسیار متفاوت هستند)، اما برای پایداری عددی، شرط بررسی می‌شود.

دوم، ساختار محاسبه به سه مرحله تفکیک شده است: محاسبه \lr{\texttt{term1}} و \lr{\texttt{term2}} (دو جمله اصلی صورت) و محاسبه \lr{\texttt{factor}} (ضریب نرمال‌سازی). این تفکیک، خوانایی کد را بالا می‌برد و امکان \lr{debug} آسان‌تر را فراهم می‌سازد. در صورتی که مقدار $Y_1^L$ منفی شود (که می‌تواند در شرایط نویز آماری رخ دهد)، \lr{\texttt{max(Y1\_L, 0)}} تضمین می‌کند که $Y_1$ نهایی غیرمنفی باشد. این برش، بازتاب‌دهنده مفهوم فیزیکی $Y_1$ به‌عنوان احتمال است که نمی‌تواند منفی باشد.

\subsection{پیاده‌سازی فرمول GLLP برای نرخ کلید}

پیاده‌سازی فرمول \lr{GLLP} برای $R_{\text{asymp}}$ مستقیم است:

\begin{equation}
	R_{\text{asymp}} = q_{\text{sift}} \cdot \left[ Q_1 \cdot (1 - H_2(e_1)) - Q_\mu^{\text{raw}} \cdot f_{\text{EC}} \cdot H_2(E_\mu) \right]
\end{equation}

در کد، این فرمول به‌صورت زیر پیاده می‌شود:

\begin{LTR}
	\begin{lstlisting}[language=Python]
		Q1 = mu * math.exp(-mu) * Y1
		R_asymp = q_sift * (Q1 * (1 - H2(e1)) - Q_mu_raw * f_ec * H2(E_mu))
		R_asymp = max(R_asymp, 0.0)
	\end{lstlisting}
\end{LTR}

نکته مهم در این پیاده‌سازی، استفاده از \lr{\texttt{Q\_mu\_raw}} (نه \lr{\texttt{Q\_mu\_sifted}}) در فرمول است. دلیل آن آن است که در فرمول \lr{GLLP}، \lr{Gain} بدون \lr{sifting} به‌کار می‌رود، چون \lr{sifting} در سطح کل فرمول (یعنی در $q_{\text{sift}}$ بیرونی) اعمال می‌شود. این یک نکته ظریف است که در پیاده‌سازی‌های نادرست به‌کررات نادیده گرفته می‌شود و می‌تواند منجر به نرخ کلید اشتباه شود.

برش \lr{\texttt{max(R\_asymp, 0.0)}} تضمین می‌کند که نرخ کلید منفی نشود. این برش در رژیم‌های با \lr{QBER} بالا (مثلاً $E_\mu > 0.11$ که آستانه بحرانی \lr{BB84} است) ضروری است، چون در این رژیم، عبارت $Q_\mu^{\text{raw}} \cdot f_{\text{EC}} \cdot H_2(E_\mu)$ بزرگ‌تر از $Q_1 \cdot (1 - H_2(e_1))$ می‌شود و $R_{\text{asymp}}$ منفی می‌گردد. در عمل، نرخ کلید منفی به این معناست که هیچ کلید امنی قابل استخراج نیست و مقدار صفر گزارش می‌شود.

\subsection{الگوریتم Lazy Import و مدیریت خطا}

پیاده‌سازی \lr{lazy import} برای \lr{\texttt{main\_optimized}} الگوی مهمی در مهندسی نرم‌افزار است. این الگو به دو دلیل به‌کار می‌رود: نخست، زمان بارگذاری ماژول را کاهش می‌دهد (در صورتی که ماژول به‌عنوان کتابخانه وارد شود و تابع \lr{\texttt{run\_absolute\_test}} فراخوانی نشود، \lr{\texttt{main\_optimized}} بارگذاری نمی‌شود). دوم، در صورت عدم دسترسی به \lr{\texttt{main\_optimized}}، خطا به‌صورت کنترل‌شده مدیریت می‌شود و تابع مقدار \lr{False} برمی‌گرداند.

\begin{LTR}
	\begin{lstlisting}[language=Python]
		try:
		import main_optimized as m
		except ImportError:
		print("Error: main_optimized not found.")
		return False
	\end{lstlisting}
\end{LTR}

این الگو، نمونه‌ای از \lr{graceful degradation} است: در صورت عدم دسترسی به یک وابستگی، سیستم به‌جای \lr{crash} کردن، رفتار قابل پیش‌بینی نشان می‌دهد. این الگو در محیط‌های \lr{CI} که ممکن است وابستگی‌ها ناقص باشند، حیاتی است.

\subsection{الگوریتم Deep Copy و جلوگیری از Mutation}

استفاده از \lr{\texttt{copy.deepcopy}} برای کپی \lr{\texttt{DEFAULT\_CONFIG}} یک الگوی حیاتی در پایتون است. در پایتون، انتساب دیکشنری (\lr{\texttt{cfg = m.DEFAULT\_CONFIG}}) فقط یک مرجع (\lr{reference}) جدید ایجاد می‌کند و تغییرات روی \lr{\texttt{cfg}} روی \lr{\texttt{DEFAULT\_CONFIG}} اصلی نیز اعمال می‌شود. این می‌تواند منجر به \lr{regression} در آزمون‌های بعدی شود، چون \lr{\texttt{DEFAULT\_CONFIG}} دیگر \emph{پیش‌فرض} نیست.

\begin{LTR}
	\begin{lstlisting}[language=Python]
		cfg = copy.deepcopy(m.DEFAULT_CONFIG)
	\end{lstlisting}
\end{LTR}

\lr{\texttt{deepcopy}} یک کپی عمیق و مستقل ایجاد می‌کند: تمام اشیای تودرتوی (دیکشنری‌ها، لیست‌ها و غیره) به‌صورت بازگشتی کپی می‌شوند و شیء جدید، کاملاً مستقل از شیء اصلی است. این الگو در فریم‌ورک‌های آزمون که در آن‌ها پیکربندی به‌کررات تغییر می‌کند، استاندارد است. هزینه \lr{\texttt{deepcopy}} از مرتبه $O(n)$ است که در اینجا $n$ تعداد کلیدهای دیکشنری است و قابل صرف‌نظر است.

\subsection{الگوریتم Relative vs Absolute Tolerance}

پیاده‌سازی تابع \lr{\texttt{check}} از الگوریتم دوگانه \lr{relative}/\lr{absolute tolerance} استفاده می‌کند:

\begin{LTR}
	\begin{lstlisting}[language=Python]
		if th_val == 0: 
		err = abs(fw_val - th_val)
		passed = err < 1e-9
		else:
		err = abs(fw_val - th_val) / abs(th_val) * 100
		passed = err <= tol
	\end{lstlisting}
\end{LTR}

این الگوریتم دو حالت را تفکیک می‌کند: در صورتی که مقدار تحلیلی صفر باشد، خطای مطلق با آستانه $10^{-9}$ مقایسه می‌شود (چون خطای نسبی در این حالت تعریف‌نشده است). در غیر این صورت، خطای نسبی به درصد محاسبه می‌شود و با تلورانس مقایسه می‌گردد. این الگوریتم، الگوی استاندارد در آزمون‌های عددی است و در کتابخانه‌هایی نظیر \lr{NumPy} (\lr{\texttt{numpy.isclose}}) و \lr{Python math.isclose} نیز به‌کار می‌رود. تفاوت در این است که در این پیاده‌سازی، کاربر می‌تواند تلورانس را به‌صورت صریح تنظیم کند، در حالی که در \lr{\texttt{math.isclose}}، تلورانس‌های پیش‌فرض $10^{-9}$ برای نسبی و $10^{-15}$ برای مطلق هستند.

\subsection{الگوریتم Comparison Pipeline و ساختار CSV}

پس از اجرای شبیه‌سازی، خروجی از فایل \lr{CSV} خوانده می‌شود:

\begin{LTR}
	\begin{lstlisting}[language=Python]
		with open("qkd_results_optimized_sweep.csv", "r") as f:
		rows = list(csv.DictReader(f))
		latest = rows[-1]
		
		fw_Q_mu = float(latest.get("detection_yield_signal", 0))
		fw_E_mu = float(latest.get("qber_signal", 0))
		fw_Y1 = float(latest.get("single_photon_yield", 0))
		fw_e1 = float(latest.get("single_photon_qber", 0))
		fw_R = float(latest.get("secure_key_rate", 0))
	\end{lstlisting}
\end{LTR}

استفاده از \lr{\texttt{csv.DictReader}} به‌جای \lr{\texttt{csv.reader}}، خوانایی را بالا می‌برد: دسترسی به ستون‌ها با نام به‌جای ایندکس. \lr{\texttt{rows[-1]}} آخرین ردیف را برمی‌گرداند که در صورتی که فقط یک نقطه فاصله آزمون شده باشد، همان ردیف آزمون است. در صورتی که فایل دارای چندین ردیف باشد (مثلاً از اجرای قبلی)، این رویکرد تضمین می‌کند که آخرین نتیجه استفاده شود. استفاده از \lr{\texttt{.get(key, 0)}} با مقدار پیش‌فرض $0$، در صورتی که ستون وجود نداشته باشد، از خطا جلوگیری می‌کند. این الگوی \lr{defensive programming} در خواندن \lr{CSV} که ممکن است ساختار آن بین نسخه‌های فریم‌ورک متغیر باشد، حیاتی است.

\subsection{محاسبه ضریب انتقال برای مقایسه}

پس از خواندن خروجی شبیه‌سازی، ضریب انتقال تحلیلی محاسبه می‌شود:

\begin{LTR}
	\begin{lstlisting}[language=Python]
		eta = channel_transmittance(dist_km, 0.2) * eta_det
		th = get_analytics(mu, nu, eta, Y0, ed, f_ec)
	\end{lstlisting}
\end{LTR}

در این خط، $\eta_{\text{channel}}$ با $\alpha_{\text{dB/km}} = 0.2$ (مقدار استاندارد فیبر تک‌حالت در $1550$ نانومتر) محاسبه می‌شود و سپس در $\eta_{\text{det}} = 0.15$ ضرب می‌گردد تا $\eta_{\text{total}}$ به دست آید. سپس تابع \lr{\texttt{get\_analytics}} با این $\eta$ فراخوانی می‌شود و مقادیر تحلیلی محاسبه می‌گردند. این الگو، تضمین می‌کند که مقایسه بین شبیه‌سازی و تحلیلی، در همان شرایط فیزیکی انجام شود.

\section{ارسال و دریافت با سایر ماژول‌ها}

\subsection{وابستگی به main\_optimized}

اصلی‌ترین وابستگیِ این ماژول، وابستگی به \lr{\texttt{main\_optimized.py}} است. این وابستگی در سه بخش نمود می‌یابد: نخست، \lr{\texttt{DEFAULT\_CONFIG}} به‌عنوان نقطهٔ آغاز پیکربندی به کار می‌رود. دوم، تابع \lr{\texttt{run\_and\_save\_csv}} برای اجرای شبیه‌سازی فراخوانی می‌شود. سوم، فایل \lr{\texttt{CSV}} با نام \lr{\texttt{qkd\_results\_optimized\_sweep.csv}} که توسط ماژول \lr{\texttt{main\_optimized}} تولید می‌شود، خوانده می‌شود.

این وابستگی، یک وابستگی \lr{tight coupling} است: تغییر در ساختار \lr{\texttt{DEFAULT\_CONFIG}} یا نام ستون‌های \lr{CSV} می‌تواند منجر به شکست این ماژول شود. برای کاهش این \lr{coupling}، در پیاده‌سازی‌های آینده می‌توان از الگوهای زیر استفاده کرد: تعریف یک \lr{schema} صریح برای ستون‌های \lr{CSV} (به‌عنوان مثال، در قالب \lr{dataclass} یا \lr{TypedDict})، یا استفاده از \lr{API} تابعی به‌جای فایل \lr{CSV} (یعنی تابع \lr{\texttt{run\_and\_save\_csv}} خروجی را به‌صورت یک شیء پایتونی برگرداند).

\subsection{دیکشنری DEFAULT\_CONFIG و Override Pattern}

\lr{\texttt{DEFAULT\_CONFIG}} یک دیکشنری تودرتوی است که در \lr{\texttt{main\_optimized.py}} تعریف شده و شامل تمام پارامترهای شبیه‌سازی است: تنظیمات شبیه‌سازی (\lr{\texttt{simulation}})، تنظیمات منبع (\lr{\texttt{source}})، تنظیمات آشکارساز (\lr{\texttt{detector}})، تنظیمات پروتکل (\lr{\texttt{protocol\_params}}) و غیره. این ماژول با \lr{\texttt{deepcopy}} یک کپی مستقل ایجاد می‌کند و سپس مقادیر مورد نظر را \lr{override} می‌نماید.

این الگوی \lr{configuration override} در فریم‌ورک‌های شبیه‌سازی استاندارد است و امکان آزمون آسان در شرایط مختلف را فراهم می‌سازد. مزیت اصلی آن آن است که فقط پارامترهای مورد نیاز برای آزمون تغییر می‌کنند و سایر پارامترها در مقدار پیش‌فرض باقی می‌مانند؛ این تضمین می‌کند که آزمون در شرایط نزدیک به واقعی انجام شود و در عین حال، تکرارپذیر باشد.

\subsection{خروجی CSV و اسکیمای ستون‌ها}

فایل \lr{CSV} که توسط \lr{\texttt{main\_optimized}} تولید می‌شود، شامل ستون‌های متعددی است که نتایج شبیه‌سازی را در بر دارد. در این ماژول، پنج ستون کلیدی خوانده می‌شوند:

\begin{itemize}
	\item \lr{\texttt{detection\_yield\_signal}}: \lr{Gain} سیگنال sifted شده ($Q_\mu^{\text{sifted}}$)
	\item \lr{\texttt{qber\_signal}}: \lr{QBER} سیگنال ($E_\mu$)
	\item \lr{\texttt{single\_photon\_yield}}: بازده تک‌فوتونی ($Y_1$)
	\item \lr{\texttt{single\_photon\_qber}}: خطای تک‌فوتونی ($e_1$)
	\item \lr{\texttt{secure\_key\_rate}}: نرخ کلید امن ($R$)
\end{itemize}

این اسکیما، در فریم‌ورک به‌عنوان \lr{contract} بین موتور شبیه‌سازی و ابزارهای تحلیل عمل می‌کند. در صورتی که ساختار \lr{CSV} تغییر کند (مثلاً ستون‌ها اضافه یا حذف شوند)، این ماژول و سایر ابزارهای تحلیل باید به‌روز شوند. برای کاهش این وابستگی، می‌توان از \lr{schema versioning} استفاده کرد: یک ستون \lr{\texttt{schema\_version}} به \lr{CSV} اضافه شود و ابزارهای تحلیل بر اساس نسخه، رفتار متفاوت داشته باشند.

\subsection{ارتباط با validate\_physics.py}

این ماژول و \lr{\texttt{validate\_physics.py}} دو لایه مکمل اعتبارسنجی را تشکیل می‌دهند. \lr{\texttt{validate\_physics.py}} بر آزمون آماری توزیع‌های منبع متمرکز است و از \lr{Z-test} برای راستی‌آزمایی \lr{Gain} و \lr{QBER} استفاده می‌کند. این ماژول (\lr{\texttt{absolute\_crypto\_validation.py}}) بر صفت کریپتوگرافیک پارامترهای کلیدی متمرکز است و از تلورانس عددی برای راستی‌آزمایی $Y_1$، $e_1$ و $R$ استفاده می‌نماید. این دو لایه، نقش‌های مکمل دارند: لایه اول (\lr{validate\_physics}) تضمین می‌کند که فیزیک پایه (توزیع \lr{Poisson} فوتون، دارک کانت و غیره) به‌درستی پیاده شده، و لایه دوم (این ماژول) تضمین می‌کند که استخراج کریپتوگرافیک (\lr{decoy-state} estimation، \lr{privacy amplification} و غیره) به‌درستی انجام شده.

\subsection{ارتباط با ماژول‌های اثبات Lim-2014 و Ma-2005}

این ماژول به‌طور غیرمستقیم با ماژول‌های اثبات امنیتی \lr{\texttt{lim2014.py}} و \lr{\texttt{ma2005.py}} در ارتباط است. از طریق \lr{\texttt{main\_optimized}}، این ماژول اثبات \lr{Lim-2014} را برای اجرای شبیه‌سازی استفاده می‌کند. با این حال، فرمول‌های تحلیلی در تابع \lr{\texttt{get\_analytics}} بر اساس حد مجانبی هستند (یعنی $N \to \infty$) که با فرمول‌های متناهی-\lr{key} در \lr{Lim-2014} متفاوت است. این تفاوت، عمدی است: آزمون کران بالای نرخ کلید را راستی‌آزمایی می‌کند و در صورتی که موتور شبیه‌سازی نتواند به کران بالا نزدیک شود، قطعاً نمی‌تواند نرخ کلید متناهی صحیح را تولید کند. این یک الگوی \lr{necessary condition testing} است که در آزمون‌های عددی به‌کررات به‌کار می‌رود.

\subsection{خروجی نهایی و استفاده در CI/CD}

خروجی نهایی این ماژول به دو شکل است: نخست، خروجی چاپی به کنسول که شامل جدول مقایسه و نتیجه \lr{PASS}/\lr{FAIL} است. دوم، کد خروج \lr{POSIX} ($0$ برای موفقیت، $1$ برای شکست) که توسط \lr{CI/CD pipeline} قابل تفسیر است. این الگو، استاندارد \lr{shell scripting} است و در ابزارهایی نظیر \lr{GitHub Actions}، \lr{GitLab CI} و \lr{Jenkins} به‌طور گسترده پشتیبانی می‌شود.

در یک \lr{CI/CD pipeline} نمونه، این ماژول به‌صورت زیر به‌کار می‌رود:

\begin{LTR}
	\begin{lstlisting}[language=bash]
		- name: Cryptographic Validation
		run: python absolute_crypto_validation.py
	\end{lstlisting}
\end{LTR}

در صورت شکست، \lr{pipeline} متوقف می‌شود و توسعه‌دهنده باید قبل از \lr{merge}، خطا را اصلاح کند. این الگو، تضمین می‌کند که هرگز کدی با خطای کریپتوگرافیک به شاخه اصلی \lr{merge} نشود. در ادبیات مهندسی نرم‌افزار علمی، این الگو به‌عنوان \lr{cryptographic regression testing} شناخته می‌شود و در پروژه‌های بالغی مانند \lr{Open Quantum Safe} و \lr{Qiskit} به‌صورت استاندارد به کار می‌رود.

\subsection{جمع‌بندی ارتباطات}

به‌طور خلاصه، این ماژول در نقش یک \lr{validator} مستقل عمل می‌کند که از طریق سه کانال با سایر بخش‌های فریم‌ورک در ارتباط است: ورودی پیکربندی از \lr{\texttt{main\_optimized.DEFAULT\_CONFIG}}، اجرای شبیه‌سازی از طریق \lr{\texttt{main\_optimized.run\_and\_save\_csv}}، و خواندن نتایج از فایل \lr{\texttt{qkd\_results\_optimized\_sweep.csv}}. خروجی آن به‌صورت کد خروج \lr{POSIX} و گزارش چاپی، برای مصرف توسط \lr{CI/CD pipeline} و توسعه‌دهنده طراحی شده است. این معماری، جداسازی واضحی بین \lr{validator} و موتور شبیه‌سازی برقرار می‌سازد و امکان توسعه مستقل هر یک را فراهم می‌نماید.
